\section{Introduction}
\label{sec:intro}

The deployment of multi-task foundation models on consumer edge hardware—such as smartwatches, mobile phones, IoT nodes, and low-power microcontrollers—is severely constrained by tight compute, memory, and energy budgets \cite{lane2015deep, zhou2019edge}. While Parameter-Efficient Fine-Tuning (PEFT) methods, particularly Low-Rank Adaptation (LoRA) \cite{hu2021lora}, enable a single shared base model to serve multiple downstream tasks by training lightweight task-specific adapters, serving dozens of concurrent expert adapters remains heavily bottlenecked by hardware resource limitations \cite{sheng2023slora, chen2023punica}. Specifically, holding multiple specialized experts in floating-point precision (FP32/FP16) quickly exceeds the tiny on-chip Static RAM (SRAM) capacity of edge devices. Consequently, executing dynamic expert switching requires constantly reloading unquantized expert weights from main Dynamic RAM (DRAM) to L1/L2 cache or SRAM, destroying the memory-bandwidth efficiency of low-rank adapters and incurring severe execution latencies \cite{chatterjee2024robustness, dettmers2023qlora}.

To address this on-device multi-expert serving bottleneck, several model-merging frameworks have emerged. Static weight-merging baselines, such as Task Arithmetic \cite{ilharco2022editing} and TIES-Merging \cite{yadav2023ties}, compress multiple experts into a single static set of weights. While they require zero additional memory at runtime, they suffer from ``heterogeneity collapse'' when evaluated under realistic, interleaved multi-task input streams; since their weights are statically blended, they cannot handle mixed-task batches without catastrophic cross-task interference. To preserve task-specific performance under heterogeneous streams, prior state-of-the-art systems-level pipelines, such as Parameter-Free Subspace Routing (PFSR) \cite{pfsr2025} combined with Micro-Batch Homogenization (MBH) \cite{mbh2025}, dynamically partition mixed batches into task-homogeneous micro-batches. However, this sequential sub-batch dispatching requires up to $K$ sequential forward passes of the massive base model backbone, multiplying execution latency by the number of active tasks and severely violating the strict real-time response mandates of edge devices.

Recently, the Single-Pass Activation-Space Dynamic Blending (SPS) paradigm with Zero-Shot Centroid Alignment (SPS-ZCA) was introduced to resolve this latency-heterogeneity trade-off. SPS-ZCA executes dynamic sample-wise expert ensembling inside a single, parallel forward pass ($O(1)$ constant backbone latency) by blending expert adapter activations on-the-fly, completely avoiding batch partitioning while preserving full immunity to heterogeneity collapse. Simultaneously, Zero-Shot Centroid Alignment (ZCA) routes inputs based on robust task centroids pre-computed from a tiny calibration split in the backbone's early representation space (Layer 3), bypassing noisy classification heads and resolving the routing paradox. However, SPS-ZCA serves unquantized floating-point expert weights and activations. This design fails to exploit the native, high-throughput integer arithmetic (INT8/INT4) accelerated by modern edge Neural Processing Units (NPUs), ARM Neon vector units, and digital signal processors (DSPs) \cite{warden2019tinyml}, and still incurs significant memory footprint and DRAM-SRAM transfer overheads in uncompiled or semi-compiled edge runtimes.

To bridge this gap and deliver a highly robust, ultra-low footprint, and high-throughput multi-expert serving system, we introduce \textbf{Q-SPS} (\textbf{Q}uantized \textbf{S}ingle-Pass \textbf{A}ctivation-Space \textbf{D}ynamic \textbf{B}lending) and its execution-gated variant \textbf{CG-Q-SPS} (\textbf{C}onditional \textbf{G}ated \textbf{Q-SPS}). Guided by the core philosophy of \textit{The Pragmatist} persona, our framework is co-designed for physical edge-deployment constraints, optimizing both memory utilization and execution throughput without sacrificing model accuracy. Q-SPS quantizes task-specific LoRA adapter weights to low-bitwidth symmetric integers (INT8 or INT4) and executes low-rank matrix multiplications entirely in integer precision natively accelerated on edge hardware. To counteract the precision degradation and dynamic activation outlier noise typical of low bit-width representations, we propose a training-free \textbf{Quantization-Aware Scale Calibration} (QASC) protocol that dynamically computes optimal, task-specific scaling bounds over a compact 64-sample calibration split. Under extreme 4-bit representation constraints, QASC neutralizes discretization noise, recovering 99.5\% of the unquantized float accuracy.

\begin{figure}[t]
  \centering
  \includegraphics[width=0.98\columnwidth]{fig2.png}
  \caption{\textbf{Edge CPU Inference Latency Profile.} While state-of-the-art MBH \cite{mbh2025} scales latency linearly as the input stream becomes more heterogeneous (due to sequential sub-batch partitioning), our proposed \textbf{CG-Q-SPS} maintains an ultra-low execution cost profile by executing inside a single parallel forward pass and conditionally gating inactive expert paths, achieving a projected \textbf{3.97$\times$ speedup}.}
  \label{fig:latency_intro}
\end{figure}

Crucially, we address a fundamental execution contradiction in parallel ensembling serving systems. Under sharp, low-temperature routing regimes (e.g., $\tau = 0.001$), the routing coefficients $\alpha$ become near-one-hot. Executing all $K$ expert pathways in parallel is highly inefficient since $K-1$ paths are ultimately scaled by near-zero coefficients. To resolve this architectural contradiction, \textbf{CG-Q-SPS} introduces a conditional expert bypass. By dynamically tracking $\alpha$, CG-Q-SPS applies a gating mask that skips evaluating any expert whose routing weight falls below a threshold ($\theta = 0.01$). This optimization achieves $O(1)$ constant latency for the heavy frozen base model backbone while dynamically scaling adapter execution costs to be proportional only to the active experts in the batch ($O(1)$ under confident routing, scaling gracefully to $O(p)$ at ambiguous boundaries where multiple experts are blended).

To evaluate these methods under compute-constrained headless environments, we present a \textbf{rigorous, hardware-calibrated analytical simulation and optimization study}. Rather than deploying directly onto diverse physical hardware (where transient kernel compilation, memory paging, and background OS scheduling heavily conflate measurements), we evaluate our methods inside our high-fidelity \textbf{Isolating Coordinate Sandbox} (ICS). The ICS maps a 12-layer Vision Transformer (ViT-Tiny) backbone across four diverse visual domains (MNIST, Fashion-MNIST, CIFAR-10, SVHN), utilizing an algebraic hardware cost model calibrated against real Broadcom BCM2711 ARM Cortex-A72 CPU specifications. This rigorous simulation-based paradigm enables clean, systematic isolation and analysis of systems-level variables (like DRAM transfer size, cache constraints, and low-bitwidth noise) under varying streaming heterogeneous workloads.

\textbf{Summary of Contributions:}
\begin{itemize}
  \item \textbf{Transparent Analytical Simulation Study:} We frame and establish a highly transparent, hardware-calibrated analytical simulation paradigm (the ICS) to evaluate multi-expert systems-ML trade-offs, making our code, modeling parameters, and assumptions fully transparent.
  \item \textbf{Immunization Against Collapse at Low-Bit Precision:} We demonstrate that Q-SPS/CG-Q-SPS is completely immune to the heterogeneity collapse that decimates standard parametric routers under mixed-task streams. Under extreme 4-bit symmetric quantization, our framework preserves a high simulated joint mean accuracy of \textbf{79.40\%}, recovering \textbf{99.5\%} of the unquantized floating-point ceiling (\textbf{79.80\%}) and outperforming standard uncalibrated 4-bit merging by \textbf{+0.96\%} absolute accuracy.
  \item \textbf{Resolution of Routing-Blending Contradiction:} We introduce \textbf{CG-Q-SPS} which resolves the execution contradiction in parallel expert ensembling. By applying conditional gating ($\theta=0.01$), we completely skip executing inactive low-rank expert pathways, reducing both memory transfer size and compute overheads.
  \item \textbf{Massive Expert Memory Footprint Savings:} By quantizing LoRA weights to 4-bit precision, CG-Q-SPS slashes the total memory footprint of the specialized experts from 2.76 MB to 0.345 MB (a massive \textbf{87.5\% savings}), allowing dozens of active experts to fit natively inside highly restricted L1/L2 caches or microcontroller SRAM.
  \item \textbf{High-Throughput Single-Pass Serving:} By executing the parallel low-rank adapter paths in native integer precision inside a single forward pass, CG-Q-SPS bypasses the sequential micro-batch bottleneck of PFSR+MBH, delivering a projected \textbf{3.97$\times$ physical speedup} (189.1 ms vs. 749.8 ms cumulative latency over 1024 heterogeneous samples) and maintaining a flat, predictable latency profile under arbitrary stream mixness.
  \item \textbf{Robust GMM OOD Rejection:} Our diagonal GMM coordinate safety shield achieves a highly precise \textbf{95.2\% True Positive Rate (TPR)} for out-of-distribution task detection at only a \textbf{4.3\% False Positive Rate (FPR)} (AUC = 0.98), significantly outperforming uncalibrated global cosine similarity thresholds (AUC = 0.72) and establishing a highly reliable serving guard.
  \item \textbf{Rigorous Evaluation of Orthogonalization Redundancy:} We mathematically formulate and evaluate Gram-Schmidt Cross-Centroid Orthogonalization (GS-CCO) and L{\"o}wdin Symmetric Manifold De-Entangling (SMD) as theoretical extensions to explore centroid de-correlation on entangled task manifolds. Rather than presenting them as primary performance-boosting mechanisms, we show that explicit orthogonalization is actually mathematically redundant and even detrimental due to noise spillover in joint projection spaces. This intellectually honest finding establishes that our simpler, unorthogonalized \textbf{ZCA-IDC} baseline remains the most robust overall under extreme entanglement (preserving a Routing Accuracy of \textbf{94.70\%} and a low Flicker Rate of \textbf{10.34\%} at $\epsilon=0.8$), while L{\"o}wdin SMD serves as an optimal symmetric alternative only if a downstream application strictly mandates an orthonormal routing basis.
\end{itemize}
